\subsection{Exact direct output formulas}
\label{subsec:new-direct-output-formulas}

The preceding proofs refer to three local inequalities.  We derive them here
from Proposition~\ref{prop:new-exact-local-set}.  This replaces the former
presentation by iterated reach maps.

For a nonsupercritical triangle with incoming reach at least $a$ and radial
reach at least $c$, let
\[
 \mathcal B_c(a)=
 \max\{b:(a,b,c)\text{ is admissible and }a+b\le1\}.
\]
This is an outgoing capacity for one triangle; it is never composed in the
proof.

For $1/2<c\le\sqrt3/2$, put
\[
 h(c)=\frac{2c}{1+\sqrt{16c^2-3}}.
\]
For $\sqrt3/2<c<1$, put
\[
 \ell(c)=\frac c2\left(1-\sqrt{4c^2-3}\right),
 \qquad r(c)=c-\ell(c).
\]
Also set
\[
 Q_-(a,c)=\frac{\sqrt{a^2-ac+c^2}}c-a,
\]
\[
 Q_+(a,c)=
 \frac{2(1-a^2)c^2}
 {2ac^2+c+
 \sqrt{(2ac^2+c)^2-4(1-c^2)(1-a^2)c^2}}.
\]

\begin{proposition}[Four direct high-radial outputs]
\label{prop:new-four-direct-outputs}
For $1/2<c\le\sqrt3/2$,
\[
 \mathcal B_c(a)=
 \begin{cases}
  1-a,&0\le a\le1-c,\\
  Q_+(a,c),&1-c<a<h(c),\\
  h(c),&a=h(c),\\
  Q_-(a,c),&h(c)<a<c,\\
  1-a,&c\le a\le1.
 \end{cases}
 \tag{6.66}
\]
For $\sqrt3/2<c<1$,
\[
 \mathcal B_c(a)=
 \begin{cases}
  1-a,&0\le a\le1-c,\\
  Q_+(a,c),&1-c<a\le\ell(c),\\
  \ell(c),&\ell(c)<a<r(c),\\
  Q_-(a,c),&r(c)\le a<c,\\
  1-a,&c\le a\le1.
 \end{cases}
 \tag{6.67}
\]
At $a=\ell(c)$, the selected value is $r(c)$; immediately to the right it is
$\ell(c)$.
\end{proposition}

\begin{proof}
Since $c>1/2$, any containing V triangle contains the radial midpoint.
Therefore $a+b\le1$, and the supercritical cell of Proposition
\ref{prop:new-exact-local-set} is absent.  On the cap $b=1-a$, the $T$-cell
inequality becomes
\[
 M(c-M)\le0,
\]
where $M=\max\{a,b\}$.  Thus the linear cap is feasible exactly when
$a\le1-c$ or $a\ge c$.

Off the linear cap, a maximal point lies on $F_L=0$ or $F_T=0$.  The selected
$L$-cell equation has the two ordered roots $\ell(c),r(c)$; the selector
$q\le0$ retains the smaller coordinate.  The ordered halves of $F_T=0$ give
$Q_+$ and $Q_-$.  Their equality boundary $a=b$ is $h(c)$ in the low-radial
range, while their $q=0$ boundaries are
$(\ell(c),r(c))$ and $(r(c),\ell(c))$ in the high range.  These conditions
produce (6.66)--(6.67).

The formal other root in the $Q_+$ quadratic violates the geometric selector
$c\le2\max\{a,b\}$ and is discarded.  At the high-range junction, the
selected $T$-cell point has value $r(c)$, whereas the adjacent selected
$L$-cell point has value $\ell(c)$; this is the genuine downward jump stated
above.
\end{proof}

\subsubsection{Low-root and threshold calculations}

Let $c=1-e$ with $0<e<1-\sqrt3/2$.  The selected low root
$\epsilon(e)=\ell(1-e)$ satisfies the quadratic
\[
 p_e(z)=z^2-(1-e)z+(1-e)^2-(1-e)^4=0.
 \tag{6.68}
\]
Substitution gives
\[
 p_e(e)=e(1-3e+4e^2-e^3)>0.
\]
Since $e<(1-e)/2$, the point $e$ lies below the smaller root.  For
$U=2e/(1-2e)$,
\[
 p_e(U)=
 -\frac{e^2}{(1-2e)^2}
 (4e^4-20e^3+37e^2-28e+3).
\]
The polynomial in parentheses is positive on the required interval, proving
\[
 e<\epsilon(e)<\frac{2e}{1-2e}.
\]
For $e\le1/8$,
\[
 p_e(2e+5e^2)=e^2(3e+4)(8e-1)\le0,
\]
which proves (6.6).

Now let $A>\epsilon(e)$ and $C\ge1-e$.  The linear branch of
Proposition~\ref{prop:new-four-direct-outputs} is no longer available.  On
the selected constant and $Q_-$ pieces, the outgoing reach is at most
$\epsilon(e)$.  On the $Q_+$ piece, the selected quadratic has its outgoing
coordinate below the same root.  Thus $B\le\epsilon(e)$, proving the direct
threshold Lemma~\ref{lem:new-direct-threshold} with no following-edge map.

\subsubsection{The strict-supercritical envelope}

Fix $0\le c<1/2$ and maximize the outgoing boundary reach over strict
supercritical pairs.  The maximum lies on the $S$-cell frontier and at the
boundary where the incoming coordinate tends to the complementary value.
Eliminating that coordinate from $F_S=0$ gives
\[
 b^2-cb+2c-1=0.
\]
The larger selected root is
\[
 M_c^{\rm sup}=\frac{c+\sqrt{c^2-8c+4}}2.
\]
The corresponding equality state lies on the boundary $a+b=1$, so it is not
strictly supercritical.  Therefore every actual strict-supercritical role has
\[
 B<M_c^{\rm sup}.
\]
This strict nonattainment is the only supercritical output used in the T3-like
and Vd1 rescuer proofs.

\subsubsection{Quarter radial envelope}

\begin{lemma}[Quarter radial envelope]
\label{lem:new-quarter-radial-envelope}
If
\[
 0\le h_0\le p,
 \qquad p+h_0<1,
\]
then
\[
 \boxed{c_{\max}(p,h_0)\le1-\frac{h_0}{4}.}
 \tag{6.69}
\]
The inequality is strict when $h_0>0$.
\end{lemma}

\begin{proof}
Only the $L$ and $T$ cells of Proposition~\ref{prop:new-exact-local-set} can
occur.  In the $L$ cell, the selected radial root is the unique root of
\[
 P_{h_0}(c)=c^4-c^2+h_0c-h_0^2.
\]
At $c_0=1-h_0/4$,
\[
 P_{h_0}(c_0)
 =\frac{h_0}{256}
 (h_0^3-16h_0^2-240h_0+128)>0
\]
for $0<h_0\le1/2$.  Since
\[
 P_{h_0}(\sqrt3/2)=-(h_0-\sqrt3/4)^2\le0
\]
and $P_{h_0}$ is increasing on the selected root interval, the selected root
is below $c_0$.

In the $T$ cell put $s=p+h_0$ and $r=\sqrt{4s^2-3}$.  Write
\[
 \gamma=\frac{p-h_0}{s}=wr,
 \qquad0\le w<1.
\]
Then
\[
 p=\frac{s(1+wr)}2,
 \qquad h_0=\frac{s(1-wr)}2,
\]
and
\[
 c_T=\frac{s(1+wr)}{1+r}.
\]
For fixed $s,r$, the function $c_T+h_0/4$ increases in $w$, so
\[
 c_T+\frac{h_0}{4}
 <\frac{s(9-r)}8.
\]
Using $4s^2=r^2+3$,
\[
256-(r^2+3)(9-r)^2
=(1-r)(r^3-17r^2+67r+13)\ge0.
\]
Hence $c_T+h_0/4<1$.  The two cells are exhausted.
\end{proof}

\subsubsection{Vd corner margins}

In the exact Vd1/Vd2 corner form, let the boundary reaches be $a,b$, let
$t>0$, and put $d=\sqrt{t^2+t+1}$.  The own-radial reach is
\[
 c_0=\frac{d-a-tb}{t+1}.
\]
If the role has positive support on the forward adjacent arm, its far endpoint
there is
\[
 u_+=\frac{d-a-tb-1}{t}.
\]
Since $d<t+1$,
\[
 c_0<1-\frac{a+tb}{t+1}
 \le1-\min\{a,b\},
 \tag{6.70}
\]
and
\[
 u_+<1-b-\frac at<1-b.
 \tag{6.71}
\]
Thus lower boundary demands $a\ge A$, $b\ge H$ imply
\[
 c_0<1-\min\{A,H\},
 \qquad
 u_+<1-H.
\]
Reflection supplies the backward adjacent-arm statements.  These two
one-line inequalities are exactly the radial separations used in the one-Vd
placement proof.
